\documentclass[10pt,a4paper]{amsart}
\usepackage[margin=19mm]{geometry}
\usepackage{amsmath,amssymb,mathtools}
\usepackage[scr=rsfs,cal=euler]{mathalfa}
\usepackage{xfrac}
\usepackage[unicode,hidelinks,bookmarks=false]{hyperref}
\newcommand{\ie}{i.e.\ }
\newcommand{\cf}{cf.\ }
\newcommand{\resp}{respectively\ }
\newcommand{\Subsection}{\subsection}
\providecommand{\nobreakdash}{\nobreak\hbox{-}}
\setlength{\emergencystretch}{2em}
\multlinegap0pt
\usepackage{amssymb}
\usepackage[shortlabels]{enumitem}
\newtheorem{theorem}{Theorem}
\def\I#1{{{\mathrm{\mathbf I}}(#1)}}
\def\IN{{\mathbb N}}
\newcommand{\Mor}{\mathsf{Mor}}
\newcommand{\Obj}{\mathsf{Obj}}
\def\calE{{\mathcal{E}}}
\newcommand{\cntSets}{{\mathsf{ODS}}}
\def\qed{\ifmmode\squareforqed\else{\unskip\nobreak\hfil
\penalty50\hskip1em\null\nobreak\hfil\squareforqed
\parfillskip=0pt\finalhyphendemerits=0\endgraf}\fi}
\def\squareforqed{\hbox{\rlap{$\sqcap$}$\sqcup$}}
\def\src{{\sigma}}
\def\tar{{\tau}}
\title{A note on computable \'{e}tale spaces}
\author{Matthew de Brecht}
\date{Source: arXiv:2604.27466v1 (CC BY 4.0)}
\begin{document}
\maketitle

\noindent\emph{Source and licence.} A note on computable \'{e}tale spaces. Version arXiv:2604.27466v1 (CC BY 4.0), distributed by arXiv under the Creative Commons Attribution 4.0 International licence, http://creativecommons.org/licenses/by/4.0/, as recorded in the arXiv licence metadata for this exact version and shown on the abstract page https://arxiv.org/abs/2604.27466v1. Full licence notice: \texttt{licenses/arxiv-2604.27466-CC-BY-4.0.txt}.\par\smallskip

\noindent\textit{Editorial scope.} Counted unit: the article's theorem theorem (source0.tex lines 363--365). The article's complete original proof of it is delivered with it (source0.tex lines 366--390, one \texttt{proof} environment of 25 lines). No mathematics is written, repaired or re-derived: the statement and the proof are the article's own lines, in the article's own notation, and no retained line is substituted or re-flowed.  No further source-local context is needed: every cross-reference of every retained line resolves inside the two delivered spans.  Nothing was omitted from any retained span, so the package carries no omission record.  No retained line contains a citation, so the wrapper carries no bibliography and no bibliography entry.  No retained line contains a figure environment, an image-inclusion command or drawing code, so no image is delivered and the package declares no asset dependency.  Editorial formatting only, in addition: an AMS \texttt{amsart} preamble that restates the article's own declarations and macros used by the retained lines, namely the environments \texttt{theorem} on separate counters, together with the standard packages those lines need.  The retained environments print in this document as Theorem 1 in order of appearance, whereas the article prints the same items with the numbering of its own full document.  The complete native source of the article as distributed in the arXiv e-print for this exact version is bundled unchanged as \texttt{sources/199474/source0.tex}.
\begin{theorem}
$\cntSets$ is equivalent to the category of overt discrete quasi-Polish spaces. Under this equivalence, the computable points of $\cntSets_\Obj$ correspond to the computably overt computably discrete effective quasi-Polish spaces, and the computable points of $\cntSets_\Mor$ correspond to the computable functions between computably overt computably discrete effective quasi-Polish spaces.
\end{theorem}

\begin{proof}
We construct a full and faithful and essentially surjective functor $\calE$ from $\cntSets$ to the category of overt discrete quasi-Polish spaces. 
Given $\equiv$ in $\cntSets_\Obj$, define $\calE_\Obj(\equiv)=\I{\equiv}$, the space of ideals of $\equiv$.  This is quasi-Polish by definition, and discrete because if $I,J\in \calE_\Obj(\equiv)$ then $I = J$ if and only if $(\exists a\in \IN)[a\in I \,\&\, a\in J]$. Finally, $\calE_\Obj(\equiv)$ is overt (relative to $\equiv$) because $[a]_\prec\not = \emptyset$ if and only if $a\equiv a$. Clearly, if $\equiv$ is a computable point then $\calE_\Obj(\equiv)$ is a computably overt computably discrete effective quasi-Polish space.

$\calE_\Mor$ maps $\langle G, \equiv_\src, \equiv_\tar \rangle$ in $\cntSets_\Mor$ to the function $f_G\colon \calE_\Obj(\equiv_\src)\to \calE_\Obj(\equiv_\tar)$ defined as
\[f_G(I) = \{b\in \IN \mid (\exists a\in I)\, G(a,b)\}\]
for each $I\in \calE_\Obj(\equiv_\src)$. It follows from the definition of ideals and the five axioms defining elements of  $\cntSets_\Mor$ that $f_G(I)$ is an ideal of $\equiv_\tar$, hence $f_G$ is a well-defined function. Clearly, $f_G$ is computable from $\langle G, \equiv_\src, \equiv_\tar \rangle$. Note that a computable point of $\cntSets_\Mor$ has computable source and target because $\src$ and $\tar$ are computable.


It is easy to see that $\calE$ is a full and faithful functor. Next, we show that each computably overt computably discrete effective quasi-Polish space is computably homeomorphic to $\calE_\Obj(\equiv)$ for some computable point $\equiv$ of $\cntSets_\Obj$. 

Assume $\prec$ is a c.e. transitive relation, that $E_\prec=\{ a\in\IN \mid [a]_\prec\not=\emptyset\}$ is c.e., and that $D_\prec \subseteq \IN\times\IN$ is c.e. and such that $I = J$ in $\I{\prec}$ if and only if there is $\langle a,b\rangle\in D_\prec$ with $a \in I$ and $b \in J$. Without loss of generality, we can assume that $\langle a,b\rangle\in D_\prec$ implies $a\in E_\prec$ and $b\in E_\prec$.

Define
\[S = \{ c\in E_\prec \mid (\exists \langle a,b\rangle\in D_\prec)[a\prec c\,\&\, b\prec c]\}.\]  
Note that for each $c\in S$ there is a unique $I\in\I{\prec}$ such that $\{I\}=[c]_\prec$, and conversely for each $I\in\I{\prec}$ there is at least one $c\in S$ with $\{I\}=[c]_\prec$. If $[a]_\prec = [b]_\prec = \{I\}$ then $a$ and $b$ must have a $\prec$-upper bound in $I$, so by defining
\[a \equiv b \iff a,b\in S \,\&\, (\exists c\in S)[a\prec c\,\&\, b\prec c],\]
we obtain that $a\equiv b$ if and only if $a,b\in S$ and $[a]_\prec = [b]_\prec=\{I\}$ for a unique $I\in\I{\prec}$. Now define $g\colon \I{\prec}\to \calE_\Obj(\equiv)$ and $h\colon \calE_\Obj(\equiv)\to\I{\prec}$ as
\begin{eqnarray*}
g(I) &=& \{ a\in I \mid a \equiv a\}\\
h(J) &=& \{ b\in S\mid (\exists a\in J)\,b\prec a \}.
\end{eqnarray*}
Then $g(I)$ is the $\equiv$-ideal composed of all $a\in S$ with $[a]_\prec = \{I\}$. Furthermore, $h(J)$ is a $\prec$-ideal because each $\equiv$-ideal is $\prec$-directed. Therefore, $g$ and $h$ are well-defined, and easily seen to be computable inverses of each other. Therefore, $\calE_\Obj(\equiv)$ is computably homeomorphic to $\I{\prec}$. Since the proof relativizes, it follows that $\calE$ is essentially surjective.
\qed
\end{proof}

\end{document}
